\documentclass[a4paper,12pt,fleqn]{article}
\usepackage[T1]{fontenc}
\usepackage{ucs}
\usepackage[utf8x]{inputenc}
\usepackage{ngerman}
\usepackage[ngerman]{babel}
\usepackage{lastpage}
\usepackage[pdftex]{color,graphicx}
\usepackage{listings}
\usepackage{pdflscape}
\usepackage{longtable}
\usepackage[inner=2cm,outer=2cm,top=1cm,bottom=1.5cm,includeheadfoot]{geometry}
\usepackage{fancyhdr}
\usepackage{url}
\usepackage{draftwatermark}
\usepackage{booktabs}
\usepackage{blindtext} 
\usepackage{framed} 
\usepackage{xcolor} 
\colorlet{shadecolor}{black} 

\usepackage{enumitem}

\SetWatermarkText{Vertraulich}
\SetWatermarkScale{4}
\SetWatermarkLightness{0.9}

\usepackage{pgfgantt}
\usepackage{amsmath,amssymb,amsfonts,amstext}
\usepackage{floatflt}
\usepackage{tikz}
\usetikzlibrary{arrows, snakes}
\usetikzlibrary{through}
\usetikzlibrary{calc}

% highlighting
\usepackage{xcolor,soul}

%---- PageLayout
\pagestyle{fancy}

\setlength{\headsep}{10mm}

\usepackage{eso-pic}

%----------------------------------------------------------------------------
% HEADER --------------------------------------------------------------------
%----------------------------------------------------------------------------
\fancyhead[R]{
  \includegraphics[width=100pt,keepaspectratio]{img/amedo2012.png}
}

\fancyhead[C]{ Wochenbericht KW 21 }

\fancyhead[L]{
  \begin{tabular}[b]{l}
  Christoph Gnip\\
  Projekt: PRPS-Evolution
  \end{tabular}
}

%Linie oben
\renewcommand{\headrulewidth}{0.5pt}
%----------------------------------------------------------------------------


%----------------------------------------------------------------------------
%----------------------------------------------------------------------------
%----------------------------------------------------------------------------
\fancyfoot[L]{Stand: \today}
\fancyfoot[C]{ EXTERN }
\fancyfoot[R]{\thepage{} von \pageref{LastPage}}

% Linie unten
\renewcommand{\footrulewidth}{0.5pt}
%----------------------------------------------------------------------------

% Import Macros  ------------------------------------------------------------
\newcommand\nn{\newline\newline}

%----------------------------------------------------------------------------
% Start the Document --------------------------------------------------------
%----------------------------------------------------------------------------
\begin{document}

\setlength{\headheight}{36pt}

\begin{titlepage}


%- the Title page --------------------------------------------------------
\begin{center}
%\vspace*{2.5cm}
{\Huge \textbf{Wochenbericht KW 21}\par}
\vspace{1cm}
{\Huge 21.5. - 26.5.2013\par}
\vspace{1cm}
{\Huge Projektwoche: 5\par}

\vspace{2cm}

\large{Erstellt durch}\\
\Large{\textbf{Christoph Gnip}}


%\vspace{4cm}
\vfill

{\normalsize Fachbereich Elektrotechnik und angewandte Naturwissenschaften\\
Westfälische Hochschule\\[2ex]Mai 2013}


\end{center}
\newpage

\end{titlepage}

%- Section 1 ----------------------------------------------------------------
\section[Allgemeines]{Allgemeines}
Ab dieser Woche wird es eine Aufteilung der Wochenberichte in eine externe- und eine interne-Version geben. Von dieser Woche an wird erwartet, dass Entwicklungen nicht mehr im Detail dargestellt werden können ohne Firmeninteressen preiszugeben. Daher wird die externe- Version nicht alle Informationen enthalten.

%- Section 2 ----------------------------------------------------------------
\section[Fortschritt]{Projektfortschritt}
In dieser Woche konnte das Modell um Wellenzahl und Phaseninformationen erweitert werden.\\
Es wurde durch analytische und numerische Lösungsmethoden der aktuelle Messaufbau, bei dem die Antennenposition frei gewählt werden kann, erfolgreich Vermessen.\\
Bis zum Ende der Woche wurde daran gearbeitet, reale Messdaten aus dem neuen Aufbau zu gewinnen. Dies ist in Ansätzen gelungen, jedoch zeigte sich bei der Verarbeitung der Daten, dass die Messwerte noch nicht sauber den Antennen zugeordnet werden können. Ausführlicher wird dies in Abschnitt~\ref{sec:Problems}

%- Section 2.2 --------------------------------------------------------------
\subsection{Modellbildung}
Das in der letzten Woche vorgestellte Modell, basierend auf der Trilateration, wurde in dieser Woche Erweitert. Das Ergebnis wird in dem internen Bericht genauer beleuchtet. Die Erweiterung führt zu einem Modell indem Wellenzahl und Phase enthalten sind. Die Phasenlage wird von dem Messaufbau ermittelt, die Wellenzahl, als final gesuchte Größe, durch num. Verfahren gefunden. Aus dem Modell werden in weiteren Schritten sinnvolle Randbedingungen (im Folgenden Constraints genannt) abgeleitet, die eine Berechnung effizient gestalten werden.\\
%
Das Modell hat im wesentlichen nun folgende Gestalt:
\begin{equation}
%\label{eq:wavenumber_trilateration_model}
0=
\underbrace{\left(
	\begin{array}{ccc}
		x_k-x_0 & y_k-y_0 & z_k-z_0 
	\end{array}
\right)}_{\textbf{A}}
\underbrace{\left(
   \begin{array}{c}
	   x-x_0\\
	   y-y_0\\
	   z-z_0
   \end{array}
\right)}_{\textbf{x}}
-
\underbrace{\left(
	\begin{array}{c}
		c_{k0}(n_0, n_k, \Theta_0, \Theta_k)
	\end{array}
\right)}_{\textbf{b}}
\end{equation}
Das Modell enthält 7 Veränderliche. Diese Menge ermöglicht es eine spätere Implementation in Shark mit den Ergebnissen des Excel Solver zu vergleichen. In dem Vektor $\mathbf{b}$ sind nun  die Konstanten $\Theta_0$-$\Theta_{N'}$ sowie die Variablen $n_0$-$n_{N'}$ enthalten.\\
%
Anhand der Koeffizientenmatrix \textbf{A}, kann eine Fehlerabschätzung erstellt werden. Da in dieser Matrix die geometrischen Gegebenheiten eingehen, kann man im Vorfeld die beste Antennenkonfiguration aus den möglichen Antennen im bestimmen. Es ergibt sich so der Vorteil, dass nicht alle Konfigurationen berechnet werden müssen um anschließend die Beste (:= die mit dem kleinsten Fehler) auszuwählen. Die Arbeiten daran werden in der nächsten Woche weitergeführt.\\
%
Das Modell wurde in Excel bereits implementiert und erwartet nun die Eingabe von Messdaten. Da dieses Ergebnis durch das Optimierungsverfahren ermittelt wurde bestätigt sich die Annahme, dass das Modell sich für eine solches Verfahren prinzipiell eignet. Um den gesamten Prozess zu vereinfachen wurde in einem weiterem Schritt einige Macros in der Excel Arbeitsmappe hinterlegt. Diese erleichtern wiederkehrende Aufgaben wie das kopieren der entsprechenden Messwerte in das betreffende Arbeitsblatt, die Automatisierung des Solvers und das erstmalige Befüllen der Solver Vorhersagen mit zufälligen Werten. \\

%- Section 2.3 --------------------------------------------------------------
\subsection{Kalibrierung}
%
Zu Beginn der Woche stand das neue Kalibrierstück zur Verfügung. Es wurden die noch zur Berechnung fehlenden Größen ermittelt und mit dem Modell aus der letzten Woche wurde die Anordnung berechnet. Es konnten alle Koordinaten der einzelnen Antennen berechnet werden.\\
%
Das Modell der letzten Woche\footnote{Wochenbericht KW 20 Anhang B} wurde in Excel implementiert und auf mehrere Arten eine Lösung gefunden. Zum Einen wurde durch analytische Herangehensweise die, für das verwendete Gleichungssystem eindeutige, Lösung bestimmt. Zum Anderen wurde auf Grundlage der gleichen Messdaten mit dem Excel Solver das Problem numerisch behandelt. Die numerischen Ergebnisse sind zu den analytischen identisch. Die Ergebnisse sind in der internen Version ausführlich in den Anhängen enthalten.\\
%
Die anschließende Verifikation mit den zur Verfügung stehenden Messmitteln zeigte eine gute Übereinstimmung der Berechnung mit dem realen Aufbau. Der Fehler der gefundenen Lösung kann mit $r_0\pm\delta_r.$ mit $\delta_r=[0.26,2.01]\%$ angegeben werden.

%- Section 3 -----------------------------------------------------------------
\section{Probleme}
\label{sec:Problems}
In dieser Woche konnten keine korrekten und aussagekräftigen Messdaten generiert werden. Zum Einen gestaltete es sich schwierig, mehrere Positionen zu finden, an denen eine ausreichende Anzahl von Antennen den Tag ließt. Zum Anderen scheint die Zuordnung der Messwerte zu den entsprechendem Antennen nicht zu stimmen. Das macht eine weitere Auswertung zur Zeit nicht möglich. Eine Lösung für diese Probleme wird in der nächsten Woche erarbeitet.

%- Appendix ------------------------------------------------------------------
%
%
%
%

\begin{appendix}
%----------------------------------------------------------------------------
%----------------------------------------------------------------------------
%----------------------------------------------------------------------------
\newpage
\begin{center}
	\huge{Anhänge}
\end{center}

\normalsize

%----------------------------------------------------------------------------
\section{Erweiterung des Trilaterationsmodells}
\label{sec:extended_trilateration_model}

{
\small
Folgende Nomenklatur und Symbole gelten für diesen Abschnitt:
%
Wie gezeigt werden konnte\footnote{Wochenbericht KW 20, Anhang B} ergibt sich für den Fall der Trilateration und der Annahme, dass vier Antennen Messwerte liefern, die Gleichung:
\begin{equation}\label{eq:final_trilateration_model}
0=
\left(
	\begin{array}{ccc}
		x_k-x_0 & y_k-y_0 & z_k-z_0 
	\end{array}
\right)
\left(
   \begin{array}{c}
	   x-x_0\\
	   y-y_0\\
	   z-z_0
   \end{array}
\right)
-
\left(
	\begin{array}{c}
		c_{kj}
	\end{array}
\right) 
\end{equation}
%
Dabei ist:
\begin{equation}\label{eq:c_kj}
	c_{kj}=\frac{1}{2}[d_{kj}^2+r_{j}^2-r_k^2]
\end{equation}
%
Ziel dieser Erweiterung ist es, einen Zusammenhang zwischen diesem Modell und der Wellenzahl zu erzeugen. Folgender Ansatz wird gewählt:
	\begin{equation}\label{eq:r_0_varrho} r(\varrho)=\frac{\lambda}{2}\left(\frac{\varrho}{2\pi}+n\right),\\\lambda=\frac{c}{f}, n:= \text{Wellenzahl}
\end{equation}
%
%
Weiterhin ist $\varrho$ die gemessene Phase, die das PRPS-System liefert und $n$ die gesuchte Wellenzahl.\\
Durch einsetzen von \eqref{eq:r_0_varrho} in \eqref{eq:c_kj}, erhalten wir:
\begin{equation}\label{eq:c_k0_extended}
	c_{kj}(\varrho_0, \varrho_k, n_0, n_k) =\frac{1}{2}\left[d_{kj}^2+\frac{\lambda^2}{4}\left(\frac{\varrho_j}{2\pi}+n_0\right)^2-\frac{\lambda^2}{4}\left(\frac{\varrho_k}{2\pi}+n_k\right)^2\right]
\end{equation}
%
Wir stellen Gleichung~\eqref{eq:c_k0_extended} um:
\begin{align}
%	
	c_{kj}(\varrho_0, \varrho_k, n_0, n_k) &= \frac{1}{2}\left\{d_{kj}^2+\frac{\lambda^2}{4}\left[\left(\frac{\varrho_j}{2\pi}\right)^2+2\frac{\varrho_j}{2\pi}n_0+n_0^2 \right.\right.\nonumber\\
	&\phantom{=}\; 
	\left.\left.-\left(\frac{\varrho_k}{2\pi}\right)^2-2\frac{\varrho_k}{2\pi}n_k-n_k^2\right]\right\}\\
%    
    &=\frac{1}{2}\left\{d_{kj}^2+\frac{\lambda^2}{4}\left[\left(\frac{\varrho_j}{2\pi}\right)^2-\left(\frac{\varrho_k}{2\pi}\right)^2 \right.\right.\nonumber\\
    &\phantom{=}\;
   	\left.\left.+2\frac{\varrho_j}{2\pi}n_0-2\frac{\varrho_k}{2\pi}n_k+n_0^2-n_k^2\right]\right\}\\
%	
	&=\frac{1}{2}d_{kj}^2+\frac{\lambda^2}{8}\left[\frac{1}{(2\pi)^2}\left(\varrho_0^2-\varrho_k^2\right) \right.\nonumber\\
	&\phantom{=}\;
	\left. +\frac{1}{\pi}\left(\varrho_0n_0-\varrho_kn_k\right)+\left(n_0^2-n_k^2\right)\right]\label{c_k0_rearragend}
\end{align}
%
Führen wir nun:
\phantomeq{c_{kj}(\varrho_0, \varrho_k, n_0, n_k)}{a_{0k} := \frac{1}{2}d_{kj}^2\nonumber}
\phantomeq{c_{kj}(\varrho_0, \varrho_k, n_0, n_k)}{a_1 := \frac{\lambda^2}{8}\nonumber}
\phantomeq{c_{kj}(\varrho_0, \varrho_k, n_0, n_k)}{a_2 := a_1\frac{1}{\pi}\nonumber}
\phantomeq{c_{kj}(\varrho_0, \varrho_k, n_0, n_k)}{a_{3kj} := a_1\frac{1}{(2\pi)^2}(\varrho_j^2-\varrho_k^2)\nonumber}
%
in Gleichung~\eqref{c_k0_rearragend} ein, erhalten die finale Form der Gleichung:
\begin{equation}
c_{kj}(\varrho_0, \varrho_k, n_0, n_k) = a_{0k}+a_1(n_0^2-n_k^2)+a_2(\varrho_0n_0-\varrho_kn_k)-a_{3kj}\label{c_k0_final_form}   
\end{equation}
%
Die Einführung der Konstanten macht zum Einen die Gleichung übersichtlicher. Zum Anderen können so, mit Blick auf eine spätere Softwareimplementation, Rechenschritte gespart werden. Das sollte sich positiv auf den späteren Berechnungsaufwand auswirken.\\
%
Im Weiteren erkennt man durch scharfes hinsehen das in Gleichung~\eqref{c_k0_final_form}, für $\varrho_k=\text{const.}$ \& $\varrho_0=\text{const.}$ gilt. Das resultiert aus der Tatsache, dass . Es ermöglicht uns zu schreiben:
\begin{equation}
c_{kj}(\varrho_0, \varrho_k, n_0, n_k) = c_{kj}(n_0, n_k)
\end{equation}
%
Im engeren Sinne einer mathematischen Funktion sollten wir die Parameter alle als Argument aufnehmen. Diese Form soll darstellen, welche Größen von Interesse sind. Im späteren Gebrauch wird diese Gleichung in der Optimierung eingesetzt werden.
Für unser Gleichungssystem aus\eqref{eq:final_trilateration_model} ergibt sich:
\begin{equation}\label{eq:wavenumber_trilateration_model}
0=
\left(
	\begin{array}{ccc}
		x_k-x_0 & y_k-y_0 & z_k-z_0 
	\end{array}
\right)
\left(
   \begin{array}{c}
	   x-x_0\\
	   y-y_0\\
	   z-z_0
   \end{array}
\right)
-
\left(
	\begin{array}{c}
		c_{kj}(n_0, n_k)
	\end{array}
\right)
\end{equation}
%
Betrachten wir nun \eqref{eq:wavenumber_trilateration_model}, wählen $N'=4$ (d.h. wir verwenden 4 Antennen) und setzen $j=0$. Wir beschreiben die Konfiguration wie folgt: Antenne 0 ist die Referenz-Antenne und Antenne 0-3 sind Messwertgeber für die Phaseninformation. 
%
\begin{equation}\label{eq:wavenumber_trilateration_model_explicit}
0=
\underbrace{\left(
	\begin{array}{ccc}
		x_1-x_0 & y_1-y_0 & z_1-z_0 \\
		x_2-x_0 & y_2-y_0 & z_2-z_0 \\
		x_3-x_0 & y_3-y_0 & z_3-z_0 
	\end{array}
\right)}_{\textbf{A}}
\underbrace{\left(
   \begin{array}{c}
	   x-x_0\\
	   y-y_0\\
	   z-z_0
   \end{array}
\right)}_{\textbf{x}}
-
\underbrace{\left(
	\begin{array}{c}
		c_{10}(n_0, n_1) \\
		c_{20}(n_0, n_2) \\
		c_{30}(n_0, n_3)
	\end{array}
\right)}_{\textbf{b}}
\end{equation}
%
\begin{equation}
\mathbf{b}=
\left(
	\begin{array}{c}
		a_{01}+a_1( n_0^2-n_1^2)+a_2(\varrho_0n_0-\varrho_1n_1)-a_{310} \\
		a_{02}+a_1(n_0^2-n_2^2)+a_2(\varrho_0n_0-\varrho_2n_2)-a_{320} \\
		a_{03}+a_1(n_0^2-n_3^2)+a_2(\varrho_0n_0-\varrho_3n_3)-a_{330}
	\end{array}
\right)
\end{equation}
%
Das Ergebnis ist ein um $\varrho$ und $n$ erweitertes Gleichungssystem. Zusätzlich enthält  es mehrere geometrische Konstanten ($a_{0k}, k=\{1,..,N-1\}$), mehrere Phasen-Konstanten ($a_{3k0}, k=\{1,..,N-1\}$), sowie zwei allgemeine ($a_1$ und $a_2$). Allgemeiner formuliert ergibt sich:
%
\begin{multline}\label{eq:final_equation}
0=
\left(
	\begin{array}{ccc}
		x_k-x_0 & y_k-y_0 & z_k-z_0 
	\end{array}
\right)
\left(
   \begin{array}{c}
	   x-x_0\\
	   y-y_0\\
	   z-z_0
   \end{array}
\right) \\
-
\left(
	\begin{array}{c}
		a_{0k}+a_1(n_0^2-n_k^2)+a_2(\varrho_0k_0-\varrho_kn_k)-a_{3kj}
	\end{array}
	\right)
\end{multline}
%
Aus Gleichung~\eqref{eq:final_equation} ist durch eine geeignete Wahl von $N'=\{4,..,N\}$ sofort ersichtlich wie viele Veränderliche sich für eine gewählte Konstellation an Antennen ergeben. Für $k$ gilt in diesem Fall $k=\{1,..,N'-1\}$.\\
%
Beispielsweise ergibt sich für das Modell aus Gleichung~\eqref{eq:final_equation} mit $N'=4$, insgesamt 7 Variablen ($\mathbf{x},n_0,n_1,n_2,n_3$) . Analog würde sich für ein Modell mit allen 8 Antennen, 11 Variablen ($\mathbf{x},n_0,..,n_7$) ergeben.
}
%----------------------------------------------------------------------------
%----------------------------------------------------------------------------
%----------------------------------------------------------------------------
\newpage
	\section{Berechnungen Messaufbau}
%
Die Form der Matrix ist Klug gewählt es lässt sich direkt sagen: Das System hat eine Lösung, da $det(\mathbf{A})\neq 0$.
	% Table generated by Excel2LaTeX from sheet 'Alle'
\begin{table}[ht]
  \caption{Dargestellt ist ein Ergebnis der Berechnung der unbekannten Antennen-Koordinaten. Die Berechnung wurde Analytisch mithilfe von Excel durchgeführt}
  \vspace{0.5cm}
    \begin{tabular}{r|rrr|r|r|r|r|rrr}
    \multicolumn{11}{l}{\textit{\textbf{Lineare Modell für die Koordinatenberechnung}}} \\
          & \multicolumn{1}{c}{\textbf{0,77}} & \multicolumn{1}{c}{\textbf{0}} & \multicolumn{1}{c|}{\textbf{0}} &       & \multicolumn{1}{c|}{\textit{\textbf{x-x\_0}}} &       & \multicolumn{1}{c|}{\textbf{c\_10}} &       &       &  \\
          & \multicolumn{1}{c}{\textbf{0}} & \multicolumn{1}{c}{\textbf{0,77}} & \multicolumn{1}{c|}{\textbf{0}} & \multicolumn{1}{c|}{*} & \multicolumn{1}{c|}{\textit{\textbf{y-y\_0}}} & =     & \multicolumn{1}{c|}{\textbf{c\_20}} &       &       &  \\
          & \multicolumn{1}{c}{\textbf{0}} & \multicolumn{1}{c}{\textbf{0}} & \multicolumn{1}{c|}{\textbf{0,8}} &       & \multicolumn{1}{c|}{\textit{\textbf{z-z\_0}}} &       & \multicolumn{1}{c|}{\textbf{c\_30}} &       &       &  \\
    \multicolumn{1}{r}{} &       &       & \multicolumn{1}{r}{} & \multicolumn{1}{r}{} & \multicolumn{1}{r}{} & \multicolumn{1}{r}{} & \multicolumn{1}{r}{} &       &       &  \\
    \multicolumn{11}{l}{\textit{\textbf{Lösung des Modells}}} \\
    \multicolumn{1}{r}{} &       &       & \multicolumn{1}{r}{\textbf{1}} & \multicolumn{1}{r}{\textbf{2}} & \multicolumn{1}{r}{\textbf{3}} & \multicolumn{1}{r}{\textbf{4}} & \multicolumn{1}{r}{\textbf{5}} & \textbf{6} & \textbf{7} & \textbf{8}\\
\cline{2-11}    \multicolumn{1}{r}{} & x-x\_0 &       & \multicolumn{1}{r}{0,479} & \multicolumn{1}{r}{-0,775} & \multicolumn{1}{r}{1,521} & \multicolumn{1}{r}{-0,924} & \multicolumn{1}{r}{1,921} & -0,556 & 1,063 & 0,454 \\
    \multicolumn{1}{r}{} & y-y\_0 & =     & \multicolumn{1}{r}{-1,012} & \multicolumn{1}{r}{-1,045} & \multicolumn{1}{r}{-1,053} & \multicolumn{1}{r}{-0,197} & \multicolumn{1}{r}{0,031} & 1,092 & 1,072 & 1,355 \\
    \multicolumn{1}{r}{} & z-z\_0 &       & \multicolumn{1}{r}{0,607} & \multicolumn{1}{r}{1,346} & \multicolumn{1}{r}{1,379} & \multicolumn{1}{r}{1,326} & \multicolumn{1}{r}{1,396} & 1,436 & 1,358 & 0,670 \\
    \multicolumn{1}{r}{} &       &       & \multicolumn{1}{r}{} & \multicolumn{1}{r}{} & \multicolumn{1}{r}{} & \multicolumn{1}{r}{} & \multicolumn{1}{r}{} &       &       &  \\
    \multicolumn{1}{r}{} & \textbf{c\_k0} & \textbf{=} & \multicolumn{8}{l}{$1/2[r_0^2-r_k^2+d_{k0}^2]$} \\
    \multicolumn{1}{r}{} &       &       & \multicolumn{1}{r}{\textbf{}} & \multicolumn{1}{r}{\textbf{}} & \multicolumn{1}{r}{\textbf{}} & \multicolumn{1}{r}{\textbf{}} & \multicolumn{1}{r}{\textbf{}} & \textbf{} & \textbf{} & \textbf{} \\
    \multicolumn{1}{r}{} & c\_10 &       & \multicolumn{1}{r}{0,369} & \multicolumn{1}{r}{-0,596} & \multicolumn{1}{r}{1,171} & \multicolumn{1}{r}{-0,711} & \multicolumn{1}{r}{1,479} & -0,428 & 0,818 & 0,349 \\
    \multicolumn{1}{r}{} & c\_20 & =     & \multicolumn{1}{r}{-0,779} & \multicolumn{1}{r}{-0,804} & \multicolumn{1}{r}{-0,811} & \multicolumn{1}{r}{-0,152} & \multicolumn{1}{r}{0,024} & 0,841 & 0,826 & 1,043 \\
    \multicolumn{1}{r}{} & c\_30 &       & \multicolumn{1}{r}{0,485} & \multicolumn{1}{r}{1,077} & \multicolumn{1}{r}{1,103} & \multicolumn{1}{r}{1,061} & \multicolumn{1}{r}{1,117} & 1,149 & 1,087 & 0,536 \\
    \multicolumn{1}{r}{} &       &       & \multicolumn{1}{r}{} & \multicolumn{1}{r}{} & \multicolumn{1}{r}{} & \multicolumn{1}{r}{} & \multicolumn{1}{r}{} &       &       &  \\
    \multicolumn{11}{l}{\textit{\textbf{Inertialkoordinaten von x\_0}}} \\
    \multicolumn{1}{r}{} & x\_0  & =     & \multicolumn{1}{r}{0} & \multicolumn{1}{r}{} & \multicolumn{1}{r}{} & \multicolumn{1}{r}{} & \multicolumn{1}{r}{} &       &       &  \\
    \multicolumn{1}{r}{} & y\_0  & =     & \multicolumn{1}{r}{0} & \multicolumn{1}{r}{} & \multicolumn{1}{r}{} & \multicolumn{1}{r}{} & \multicolumn{1}{r}{} &       &       &  \\
    \multicolumn{1}{r}{} & z\_0  & =     & \multicolumn{1}{r}{0} & \multicolumn{1}{r}{} & \multicolumn{1}{r}{} & \multicolumn{1}{r}{} & \multicolumn{1}{r}{} &       &       &  \\
    \multicolumn{1}{r}{} &       &       & \multicolumn{1}{r}{} & \multicolumn{1}{r}{} & \multicolumn{1}{r}{} & \multicolumn{1}{r}{} & \multicolumn{1}{r}{} &       &       &  \\
    \multicolumn{11}{l}{\textit{\textbf{Ermittelte Koordinaten}}} \\
    \multicolumn{1}{r}{\textit{\textbf{}}} &       &       & \multicolumn{1}{r}{\textbf{1}} & \multicolumn{1}{r}{\textbf{2}} & \multicolumn{1}{r}{\textbf{3}} & \multicolumn{1}{r}{\textbf{4}} & \multicolumn{1}{r}{\textbf{5}} & \textbf{6} & \textbf{7} & \textbf{8}\\
\cline{2-11}    \multicolumn{1}{r}{} & \textbf{x\_i} & \textbf{=} & \multicolumn{1}{r}{0,479} & \multicolumn{1}{r}{-0,775} & \multicolumn{1}{r}{1,521} & \multicolumn{1}{r}{-0,924} & \multicolumn{1}{r}{1,921} & -0,556 & 1,063 & 0,454\\
    \multicolumn{1}{r}{} & \textbf{y\_i} & \textbf{=} & \multicolumn{1}{r}{-1,012} & \multicolumn{1}{r}{-1,045} & \multicolumn{1}{r}{-1,053} & \multicolumn{1}{r}{-0,197} & \multicolumn{1}{r}{0,031} & 1,092 & 1,072 & 1,355 \\
    \multicolumn{1}{r}{} & \textbf{z\_i} & \textbf{=} & \multicolumn{1}{r}{0,607} & \multicolumn{1}{r}{1,346} & \multicolumn{1}{r}{1,379} & \multicolumn{1}{r}{1,326} & \multicolumn{1}{r}{1,396} & 1,436 & 1,358 & 0,670 \\
    \multicolumn{1}{r}{} &       &       & \multicolumn{1}{r}{} & \multicolumn{1}{r}{} & \multicolumn{1}{r}{} & \multicolumn{1}{r}{} & \multicolumn{1}{r}{} &       &       &  \\
    \multicolumn{11}{l}{\textit{\textbf{Fehlerbetrachtung}}} \\
    \multicolumn{1}{r}{} &       &       & \multicolumn{1}{r}{\textbf{1}} & \multicolumn{1}{r}{\textbf{2}} & \multicolumn{1}{r}{\textbf{3}} & \multicolumn{1}{r}{\textbf{4}} & \multicolumn{1}{r}{\textbf{5}} & \textbf{6} & \textbf{7} & \textbf{8}\\
\cline{2-11}    \multicolumn{1}{r}{} & \multicolumn{2}{r}{gemessen [m]} & \multicolumn{1}{r}{1,259} & \multicolumn{1}{r}{1,894} & \multicolumn{1}{r}{2,334} & \multicolumn{1}{r}{1,661} & \multicolumn{1}{r}{2,399} & 1,851 & 2,055 & 1,574\\
    \multicolumn{1}{r}{} & \multicolumn{2}{r}{berechnet [m]} & \multicolumn{1}{r}{1,274} & \multicolumn{1}{r}{1,872} & \multicolumn{1}{r}{2,307} & \multicolumn{1}{r}{1,628} & \multicolumn{1}{r}{2,375} & 1,887 & 2,031 & 1,578 \\
    \multicolumn{1}{r}{} & \multicolumn{2}{r}{Fehler ABS [m]} & \multicolumn{1}{r}{0,015} & \multicolumn{1}{r}{0,022} & \multicolumn{1}{r}{0,027} & \multicolumn{1}{r}{0,033} & \multicolumn{1}{r}{0,024} & 0,036 & 0,024 & 0,004 \\
    \multicolumn{1}{r}{} & \multicolumn{2}{r}{\textbf{Fehler \%}} & \multicolumn{1}{r}{\textbf{1,14\%}} & \multicolumn{1}{r}{\textbf{1,19\%}} & \multicolumn{1}{r}{\textbf{1,15\%}} & \multicolumn{1}{r}{\textbf{2,01\%}} & \multicolumn{1}{r}{\textbf{1,01\%}} & \textbf{1,93\%} & \textbf{1,19\%} & \textbf{0,26\%} \\
    \end{tabular}%
  \label{tab:addlabel}%
\end{table}%

	\label{sec:coordinate_Measurements}
%----------------------------------------------------------------------------
%----------------------------------------------------------------------------
%----------------------------------------------------------------------------
\newpage
	\section{Berechnungen Messaufbau}
%
	% Table generated by Excel2LaTeX from sheet 'numerisch'
\begin{table}[h]
  \centering
  \caption{Hier gezeigt wird das Ergebnis der Berechnung der Antennenposition. Die Datengrundlage ist dieselbe wie die in Anhang~\ref{sec:coordinate_Measurements}. Die Ergebnisse sind identisch und wurden mit dem Excel Solver bestimmt. Dabei wurde die quadratische Fehlersumme minimiert. ($min(\Sigma\epsilon^2)$).}
    \vspace{0.5cm}
    \begin{tabular}{r|rrr|r|r|r|r|rrr}
    \multicolumn{11}{l}{\textit{\textbf{Lineares Modell für die Koordinatenberechnung}}} \\
          & \multicolumn{1}{c}{{0,77}} & \multicolumn{1}{c}{{0}} & \multicolumn{1}{c|}{{0}} &       & \multicolumn{1}{c|}{${x-x_0}$} &       & \multicolumn{1}{c|}{${c_{10}}$} &       &       &  \\
          & \multicolumn{1}{c}{{0}} & \multicolumn{1}{c}{{0,77}} & \multicolumn{1}{c|}{{0}} & \multicolumn{1}{c|}{*} & \multicolumn{1}{c|}{${y-y_0}$} & \multicolumn{1}{c|}{=} & \multicolumn{1}{c|}{${c_{20}}$} &       &       &  \\
          & \multicolumn{1}{c}{{0}} & \multicolumn{1}{c}{{0}} & \multicolumn{1}{c|}{{0,8}} &       & \multicolumn{1}{c|}{${z-z_0}$} &       & \multicolumn{1}{c|}{${c_{30}}$} &       &       &  \\
    \multicolumn{1}{r}{} & \multicolumn{1}{c}{\textbf{}} & \multicolumn{1}{c}{\textbf{}} & \multicolumn{1}{c}{\textbf{}} & \multicolumn{1}{r}{} & \multicolumn{1}{c}{\textit{\textbf{}}} & \multicolumn{1}{r}{} & \multicolumn{1}{c}{\textbf{}} &       &       &  \\
    \multicolumn{1}{r}{} & \multicolumn{3}{c}{\textbf{A}} & \multicolumn{1}{c}{*} & \multicolumn{1}{c}{\textit{\textbf{x}}} & \multicolumn{1}{c}{=} & \multicolumn{1}{c}{\textbf{b}} & \textbf{} &       &  \\
    \multicolumn{11}{r}{} \\
    \multicolumn{11}{l}{\textit{\textbf{Vorhersage vom Solver}}} \\
    \multicolumn{1}{r}{} &       &       & \multicolumn{1}{r}{\textbf{1}} & \multicolumn{1}{r}{\textbf{2}} & \multicolumn{1}{r}{\textbf{3}} & \multicolumn{1}{r}{\textbf{4}} & \multicolumn{1}{r}{\textbf{5}} & \textbf{6} & \textbf{7} & \textbf{8} \\
\cline{2-11}    \multicolumn{1}{r}{} & $c_{10}'$ &       & \multicolumn{1}{r}{0,369} & \multicolumn{1}{r}{-0,596} & \multicolumn{1}{r}{1,171} & \multicolumn{1}{r}{-0,711} & \multicolumn{1}{r}{1,479} & -0,428 & 0,818 & 0,349 \\
    \multicolumn{1}{r}{} & $c_{20}'$ &       & \multicolumn{1}{r}{-0,779} & \multicolumn{1}{r}{-0,804} & \multicolumn{1}{r}{-0,811} & \multicolumn{1}{r}{-0,152} & \multicolumn{1}{r}{0,024} & 0,841 & 0,826 & 1,043 \\
    \multicolumn{1}{r}{} & {$c_{30}'$} &       & \multicolumn{1}{r}{0,485} & \multicolumn{1}{r}{1,077} & \multicolumn{1}{r}{1,103} & \multicolumn{1}{r}{1,061} & \multicolumn{1}{r}{1,117} & 1,149 & 1,087 & 0,536 \\
    \multicolumn{1}{r}{} & $min(\Sigma\epsilon^2)$ &       & \multicolumn{1}{r}{0,000} & \multicolumn{1}{r}{0,000} & \multicolumn{1}{r}{0,000} & \multicolumn{1}{r}{0,000} & \multicolumn{1}{r}{0,000} & 0,000 & 0,000 & 0,000 \\
    \multicolumn{1}{r}{} &       &       & \multicolumn{1}{r}{} & \multicolumn{1}{r}{} & \multicolumn{1}{r}{} & \multicolumn{1}{r}{} & \multicolumn{1}{r}{} &       &       &  \\
    \multicolumn{11}{l}{\textit{\textbf{Lösung des Modells}}} \\
    \multicolumn{1}{r}{} &       &       & \multicolumn{1}{r}{\textbf{1}} & \multicolumn{1}{r}{\textbf{2}} & \multicolumn{1}{r}{\textbf{3}} & \multicolumn{1}{r}{\textbf{4}} & \multicolumn{1}{r}{\textbf{5}} & \textbf{6} & \textbf{7} & \textbf{8} \\
\cline{2-11}    \multicolumn{1}{r}{} &       &       & \multicolumn{1}{r}{0,479} & \multicolumn{1}{r}{-0,775} & \multicolumn{1}{r}{1,521} & \multicolumn{1}{r}{-0,924} & \multicolumn{1}{r}{1,921} & -0,556 & 1,063 & 0,454\\
    \multicolumn{1}{r}{} & $\mathbf{x_k}$ & =     & \multicolumn{1}{r}{-1,012} & \multicolumn{1}{r}{-1,045} & \multicolumn{1}{r}{-1,053} & \multicolumn{1}{r}{-0,197} & \multicolumn{1}{r}{0,031} & 1,092 & 1,072 & 1,355 \\
    \multicolumn{1}{r}{} &       &       & \multicolumn{1}{r}{0,607} & \multicolumn{1}{r}{1,346} & \multicolumn{1}{r}{1,379} & \multicolumn{1}{r}{1,326} & \multicolumn{1}{r}{1,396} & 1,436 & 1,358 & 0,670 \\
    \multicolumn{1}{r}{} &       &       & \multicolumn{1}{r}{} & \multicolumn{1}{r}{} & \multicolumn{1}{r}{} & \multicolumn{1}{r}{} & \multicolumn{1}{r}{} &       &       &  \\
    \multicolumn{1}{r}{} &       &       & \multicolumn{1}{r}{0,369} & \multicolumn{1}{r}{-0,596} & \multicolumn{1}{r}{1,171} & \multicolumn{1}{r}{-0,711} & \multicolumn{1}{r}{1,479} & -0,428 & 0,818 & 0,349 \\
    \multicolumn{1}{r}{} & $\mathbf{b_k}$ & =     & \multicolumn{1}{r}{-0,779} & \multicolumn{1}{r}{-0,804} & \multicolumn{1}{r}{-0,811} & \multicolumn{1}{r}{-0,152} & \multicolumn{1}{r}{0,024} & 0,841 & 0,826 & 1,043 \\
    \multicolumn{1}{r}{} &       &       & \multicolumn{1}{r}{0,485} & \multicolumn{1}{r}{1,077} & \multicolumn{1}{r}{1,103} & \multicolumn{1}{r}{1,061} & \multicolumn{1}{r}{1,117} & 1,149 & 1,087 & 0,536 \\
    \multicolumn{1}{r}{} &       &       & \multicolumn{1}{r}{} & \multicolumn{1}{r}{} & \multicolumn{1}{r}{} & \multicolumn{1}{r}{} & \multicolumn{1}{r}{} &       &       &  \\
    \multicolumn{11}{l}{\textit{\textbf{Inertialkoordinaten von x\_0}}} \\
    \multicolumn{1}{r}{} & x\_0  & =     & \multicolumn{1}{r}{0} & \multicolumn{1}{r}{} & \multicolumn{1}{r}{} & \multicolumn{1}{r}{} & \multicolumn{1}{r}{} &       &       &  \\
    \multicolumn{1}{r}{} & y\_0  & =     & \multicolumn{1}{r}{0} & \multicolumn{1}{r}{} & \multicolumn{1}{r}{} & \multicolumn{1}{r}{} & \multicolumn{1}{r}{} &       &       &  \\
    \multicolumn{1}{r}{} & z\_0  & =     & \multicolumn{1}{r}{0} & \multicolumn{1}{r}{} & \multicolumn{1}{r}{} & \multicolumn{1}{r}{} & \multicolumn{1}{r}{} &       &       &  \\
    \multicolumn{1}{r}{} &       &       & \multicolumn{1}{r}{} & \multicolumn{1}{r}{} & \multicolumn{1}{r}{} & \multicolumn{1}{r}{} & \multicolumn{1}{r}{} &       &       &  \\
    \multicolumn{11}{l}{\textit{\textbf{Ermittelte Koordinaten}}} \\
    \multicolumn{1}{r}{\textit{\textbf{}}} &       &       & \multicolumn{1}{r}{\textbf{1}} & \multicolumn{1}{r}{\textbf{2}} & \multicolumn{1}{r}{\textbf{3}} & \multicolumn{1}{r}{\textbf{4}} & \multicolumn{1}{r}{\textbf{5}} & \textbf{6} & \textbf{7} & \textbf{8} \\
\cline{2-11}    \multicolumn{1}{r}{} & ${x_k}$ & \textbf{=} & \multicolumn{1}{r}{0,479} & \multicolumn{1}{r}{-0,775} & \multicolumn{1}{r}{1,521} & \multicolumn{1}{r}{-0,924} & \multicolumn{1}{r}{1,921} & -0,556 & 1,063 & 0,454 \\
    \multicolumn{1}{r}{} & ${y_k}$ & \textbf{=} & \multicolumn{1}{r}{-1,012} & \multicolumn{1}{r}{-1,045} & \multicolumn{1}{r}{-1,053} & \multicolumn{1}{r}{-0,197} & \multicolumn{1}{r}{0,031} & 1,092 & 1,072 & 1,355 \\
    \multicolumn{1}{r}{} & ${z_k}$ & \textbf{=} & \multicolumn{1}{r}{0,607} & \multicolumn{1}{r}{1,346} & \multicolumn{1}{r}{1,379} & \multicolumn{1}{r}{1,326} & \multicolumn{1}{r}{1,396} & 1,436 & 1,358 & 0,670 \\
    \multicolumn{1}{r}{} &       &       & \multicolumn{1}{r}{} & \multicolumn{1}{r}{} & \multicolumn{1}{r}{} & \multicolumn{1}{r}{} & \multicolumn{1}{r}{} &       &       &  \\
    \multicolumn{11}{l}{\textit{\textbf{Fehlerbetrachtung}}} \\
    \multicolumn{1}{r}{} &       &       & \multicolumn{1}{r}{\textbf{1}} & \multicolumn{1}{r}{\textbf{2}} & \multicolumn{1}{r}{\textbf{3}} & \multicolumn{1}{r}{\textbf{4}} & \multicolumn{1}{r}{\textbf{5}} & \textbf{6} & \textbf{7} & \textbf{8} \\
\cline{2-11}    \multicolumn{1}{r}{} & \multicolumn{2}{r}{gemessen [m]} & \multicolumn{1}{r}{1,259} & \multicolumn{1}{r}{1,894} & \multicolumn{1}{r}{2,334} & \multicolumn{1}{r}{1,661} & \multicolumn{1}{r}{2,399} & 1,851 & 2,055 & 1,574 \\
    \multicolumn{1}{r}{} & \multicolumn{2}{r}{berechnet [m]} & \multicolumn{1}{r}{1,274} & \multicolumn{1}{r}{1,872} & \multicolumn{1}{r}{2,307} & \multicolumn{1}{r}{1,628} & \multicolumn{1}{r}{2,375} & 1,887 & 2,031 & 1,578 \\
    \multicolumn{1}{r}{} & \multicolumn{2}{r}{Fehler ABS [m]} & \multicolumn{1}{r}{0,015} & \multicolumn{1}{r}{0,022} & \multicolumn{1}{r}{0,027} & \multicolumn{1}{r}{0,033} & \multicolumn{1}{r}{0,024} & 0,036 & 0,024 & 0,004 \\
    \multicolumn{1}{r}{} & \multicolumn{2}{r}{\textbf{Fehler \%}} & \multicolumn{1}{r}{\textbf{1,14\%}} & \multicolumn{1}{r}{\textbf{1,19\%}} & \multicolumn{1}{r}{\textbf{1,15\%}} & \multicolumn{1}{r}{\textbf{2,01\%}} & \multicolumn{1}{r}{\textbf{1,01\%}} & \textbf{1,93\%} & \textbf{1,19\%} & \textbf{0,26\%} \\
    \end{tabular}%
  \label{tab:addlabel}%
\end{table}%

	\label{sec:coordinate_Measurements_Numerical}
%----------------------------------------------------------------------------
%----------------------------------------------------------------------------
%----------------------------------------------------------------------------
\newpage
\begin{landscape}
	\section{Projektlaufplan KW 21}
	\label{sec:projectplan}
	\scalebox{.75}{
		\begin{ganttchart}[vgrid={draw=none,*1{gray, dashed}},
				hgrid=true,
				today=12,
				title height=1,
				y unit title=0.6cm,
				y unit chart=0.8cm,
				group right shift=0,
				group top shift=.3,
				group height=.3,
				milestone width=.8,
				group peaks={}{}{.2},
				incomplete/.style={fill=black!15}, %
				bar/.style={fill=white}, %
				today label={Heute},
				today rule/.style={dashed, thick}]{44}


\gantttitle{\textbf{2013}}{44} \\
\gantttitlelist{16,...,37}{2} \\
%-------------------------------------------------------------
\ganttgroup{Projekt Evaluation}{3}{14} \\
\ganttbar[progress=100, progress label font=\small\color{black!75},
	progress label anchor/.style={right=4pt}]{Installation der Umgebungen}{3}{6} \\
	
\ganttbar[progress=100, progress label font=\small\color{black!75},
	progress label anchor/.style={right=4pt},
	bar label font=\normalsize\color{black},
	name=rech]{Recherche}{3}{7} \\
	
\ganttmilestone[name=ms1]{Vorstellung der Ergebnisse}{7} \\
	
\ganttbar[progress=90, progress label font=\small\color{black!75},
	progress label anchor/.style={right=4pt},
	bar label font=\normalsize\color{black},
	name=pflichten]
	{Pflichtenheft}{5}{8} \\
	
\ganttmilestone[name=ms2]{Pflichtenheft fertig}{8} \\

\ganttbar[progress=70, progress label font=\small\color{black!75},
	progress label anchor/.style={right=4pt},
	bar label font=\normalsize\color{black},
	name=bNumVerf]
	{Einarbeitung num. Verfahren}{5}{12} \\

\ganttbar[progress=20, progress label font=\small\color{black!75},
	progress label anchor/.style={right=34pt},
	bar label font=\normalsize\color{black},
	name=bCMAES]
	{speziell CMA-ES}{7}{10} \\

\ganttmilestone[name=ms3]{Beurteilung num. Verfahren}{12} \\

\ganttlinkedbar[progress=10, progress label font=\small\color{black!75},
	progress label anchor/.style={right=34pt},
	bar label font=\normalsize\color{black}]
	{Shark Einarbeitung}{13}{14} \\

\ganttlinkedmilestone[name=ms7]{Abschluss Evaluation}{14} \\
	
%-------------------------------------------------------------
\ganttgroup{Erstellung Prototyp}{15}{26} \\
\ganttgroup{(optional)}{15}{18} \\
\ganttbar[progress=25, progress label font=\small\color{black!75},
	progress label anchor/.style={right=4pt},
	bar label font=\normalsize\color{black}]
	{(Entwurf digi. Filter)}{15}{15} \\

\ganttlinkedbar[progress=10, progress label font=\small\color{black!75},
	progress label anchor/.style={right=4pt},
	bar label font=\normalsize\color{black},
	name=bImpFPGA]
	{(Implementation FPGA)}{16}{18} \\

\ganttmilestone[name=ms4]{(Verifikation dig. Filter)}{18} \\
	
\ganttbar[progress=0, progress label font=\small\color{black!75},
	progress label anchor/.style={right=4pt},
	bar label font=\normalsize\color{black},
	name=bImplAlgo]
	{Implementation Algorithmus}{15}{26} \\

\ganttlinkedmilestone[name=ms5]{Implementation Done}{26} \\

%-------------------------------------------------------------
\ganttgroup{Verifikation}{27}{34} \\
\ganttbar[progress=0, progress label font=\small\color{black!75},
	progress label anchor/.style={right=4pt},
	bar label font=\normalsize\color{black},
	name=bVerf]
	{Durchf\"uhrung Verifikation}{27}{34} \\

\ganttlinkedmilestone[name=ms6]{Verifikation Done}{34} \\

%-------------------------------------------------------------
\ganttgroup{Projektdokumentation}{35}{42} \\

\ganttbar[progress=0, progress label font=\small\color{black!75},
	progress label anchor/.style={right=4pt},
	bar label font=\normalsize\color{black},
	name=thesis]
	{Thesis schreiben}{35}{42} \\
	
\ganttmilestone[name=msthesis,milestone label font=\color{red}, 
	milestone/.style={fill=red}]{Abgabe}{42}

%\ganttlink{ms7}{bImplAlgo}
\ganttlink{bImpFPGA}{ms4}
\ganttlink{bNumVerf}{ms3}
\ganttlink{bCMAES}{ms3}
\ganttlink{rech}{ms1}
\ganttlink{pflichten}{ms2}
\ganttlink{thesis}{msthesis}

%
%\gantttitle{17}{1}
%\gantttitle{18}{1}
%\gantttitle{19}{1}
%\gantttitle{20}{1}
%\gantttitle{21}{1}
%\gantttitle{KW22}{1}
%\gantttitle{KW23}{1}
%\gantttitle{KW24}{1}
%\gantttitle{2011}{12} \\
%\ganttbar
%[progress=100, progress label font=\small\color{black!75},
%progress label anchor/.style={right=4pt},
%bar label font=\normalsize\color{black},
%name=pp]
%{Preliminary Project}{1}{4} \\
%\ganttset{progress label text={}, link/.style={black, -to}}
%\ganttgroup{Objective 1}{5}{16} \\
%\ganttbar[progress=4, name=T1A]{Task A}{5}{10} \\
%\ganttlinkedbar[progress=0]{Task B}{11}{16} \\
%\ganttgroup{Objective 2}{5}{16} \\
%\ganttbar[progress=15, name=T2A]{Task A}{5}{13} \\
%\ganttlinkedbar[progress=0]{Task B}{14}{16} \\
%\ganttgroup{Objective 3}{9}{12} \\
%\ganttbar[progress=0]{Task A}{9}{12}
%\ganttset{link/.style={OliveGreen}}
%\ganttlink[link mid=.4]{pp}{T1A}
%\ganttlink[link mid=.159]{pp}{T2A}

%	[canvas/
%	.style={fill=none, draw=black!5, line width=.75pt},
%		hgrid style/.style={draw=black!5, line width=.75pt},
%		vgrid={*1{draw=black!5, line width=.75pt}},
%		today=7.1,
%		today rule/.style={draw=black!64,
%		dash pattern=on 3.5pt off 4.5pt, line width=1.5pt},
%		today label={\small\bfseries TODAY},
%		title/.style={draw=none, fill=none},
%		title label font=\bfseries\footnotesize,
%		title label anchor/.style={below=7pt},
%		include title in canvas=false,
%		bar label font=\mdseries\small\color{black!70},
%		bar label anchor/.style={left=2cm},
%		bar/.style={draw=none, fill=white},
%		bar incomplete/.style={fill=barblue},
%		progress label font=\mdseries\footnotesize\color{black!70},
%		group incomplete/.style={fill=groupblue},
%		group left shift=0,
%		group right shift=0,
%		group height=.5,
%		group peaks={0}{}{},
%		group label anchor/.style={left=.6cm},
%		link/.style={-latex, line width=1.5pt},
%		link label font=\scriptsize\bfseries,
%		link label anchor/.style={below left=-2pt and 0pt}
%	]{13}
	
%	\gantttitle[title label anchor/.style={below left=7pt and -3pt}]%
%	{WEEKS:\quad1}{1}
%	\gantttitlelist{1,...,13}{1} \\
%	\ganttgroup[progress=57, progress label font=\bfseries\small]%
%	{WBS 1 Summary Element 1}{1}{10} \\
%	\ganttbar[progress=75, name=WBS1A]%
%	{\textbf{WBS 1.1} Activity A}{1}{8} \\
%	\ganttbar[progress=67, name=WBS1B]%
%	{\textbf{WBS 1.2} Activity B}{1}{3} \\
%	\ganttbar[progress=50, name=WBS1C]%
%	{\textbf{WBS 1.3} Activity C}{4}{10} \\
%	\ganttbar[progress=0, name=WBS1D]%
%	{\textbf{WBS 1.4} Activity D}{4}{10} \\[grid]
%	\ganttgroup[progress=0, progress label font=\bfseries\small]%
%	{WBS 2 Summary Element 2}{4}{10} \\
%	\ganttbar[progress=0]{\textbf{WBS 2.1} Activity E}{4}{5} \\
%	\ganttbar[progress=0]{\textbf{WBS 2.2} Activity F}{6}{8} \\
%	\ganttbar[progress=0]{\textbf{WBS 2.3} Activity G}{9}{10}
%	\ganttlink[link type=s-s]{WBS1A}{WBS1B}
%	\ganttlink[link type=f-s]{WBS1B}{WBS1C}
%	\ganttlink[link type=f-f, link label anchor/.style={left}]{WBS1C}{WBS1D}
%	
	\end{ganttchart}
		}
\end{landscape}

%----------------------------------------------------------------------------

\end{appendix}


\newpage
%- Bibliography --------------------------------------------------------------
\bibliographystyle{ieeetr}
\bibliography{../bib/mathesis_collection1}

\end{document}